% =============================================================================
%  RÉFÉRENCES BIBLIOGRAPHIQUES, norme APA (6e édition)
% =============================================================================
\cleardoublepage
\entete{Références bibliographiques}
\chapter*{RÉFÉRENCES BIBLIOGRAPHIQUES}
\addcontentsline{toc}{chapter}{RÉFÉRENCES BIBLIOGRAPHIQUES}

% Environnement à retrait négatif (« hanging indent »), conforme à l'APA.
\newenvironment{apalist}
  {\begin{list}{}{%
     \setlength{\leftmargin}{1.27cm}%
     \setlength{\itemindent}{-1.27cm}%
     \setlength{\listparindent}{0pt}%
     \setlength{\itemsep}{0.55em}%
     \setlength{\parsep}{0pt}%
     \setlength{\topsep}{0.5em}}\sloppy}
  {\end{list}}

Les références ci-après suivent la norme de l'\emph{American Psychological
Association}, classées par ordre alphabétique d'auteurs au sein de chaque
catégorie de documents.

\section*{Ouvrages}
\addcontentsline{toc}{section}{Ouvrages}

\begin{apalist}
\item Bass, L., Weber, I., \& Zhu, L. (2015). \emph{DevOps: A software
      architect's perspective}. Upper Saddle River, NJ: Addison-Wesley.

\item Humble, J., \& Farley, D. (2010). \emph{Continuous delivery: Reliable
      software releases through build, test, and deployment automation}. Boston,
      MA: Addison-Wesley.

\item Kim, G., Humble, J., Debois, P., \& Willis, J. (2016). \emph{The DevOps
      handbook: How to create world-class agility, reliability, and security in
      technology organizations}. Portland, OR: IT Revolution Press.

\item Sarr, F., \& Savoy, B. (2018). \emph{Restituer le patrimoine africain :
      vers une nouvelle éthique relationnelle}. Paris: Philippe Rey / Seuil.
\end{apalist}

\section*{Articles de périodiques et actes de colloque}
\addcontentsline{toc}{section}{Articles de périodiques et actes de colloque}

\begin{apalist}
\item Keumoe, G. D., \& Blot, C. (2023). Transformer les pratiques muséales au
      Cameroun grâce à la visualisation 3D. \emph{Humanités numériques, 7}.
      Retrieved from \url{https://journals.openedition.org/revuehn/3315}

\item Hevner, A. R., March, S. T., Park, J., \& Ram, S. (2004). Design science in
      information systems research. \emph{MIS Quarterly, 28}(1), 75--105.
      doi:10.2307/25148625

\item Ji, Z., Lee, N., Frieske, R., Yu, T., Su, D., Xu, Y., \ldots{} Fung, P.
      (2023). Survey of hallucination in natural language generation. \emph{ACM
      Computing Surveys, 55}(12), 1--38. doi:10.1145/3571730

\item Lewis, P., Perez, E., Piktus, A., Petroni, F., Karpukhin, V., Goyal, N.,
      \ldots{} Kiela, D. (2020). Retrieval-augmented generation for
      knowledge-intensive NLP tasks. \emph{Advances in Neural Information
      Processing Systems, 33}, 9459--9474.

\item Malkov, Y. A., \& Yashunin, D. A. (2018). Efficient and robust approximate
      nearest neighbor search using hierarchical navigable small world graphs.
      \emph{IEEE Transactions on Pattern Analysis and Machine Intelligence,
      42}(4), 824--836. doi:10.1109/TPAMI.2018.2889473

\item Milgram, P., \& Kishino, F. (1994). A taxonomy of mixed reality visual
      displays. \emph{IEICE Transactions on Information and Systems, E77-D}(12),
      1321--1329.

\item Slater, M., \& Wilbur, S. (1997). A framework for immersive virtual
      environments (FIVE): Speculations on the role of presence in virtual
      environments. \emph{Presence: Teleoperators and Virtual Environments,
      6}(6), 603--616. doi:10.1162/pres.1997.6.6.603

\item Vaswani, A., Shazeer, N., Parmar, N., Uszkoreit, J., Jones, L., Gomez,
      A. N., \ldots{} Polosukhin, I. (2017). Attention is all you need.
      \emph{Advances in Neural Information Processing Systems, 30}, 5998--6008.
\end{apalist}

\section*{Rapports et documents institutionnels}
\addcontentsline{toc}{section}{Rapports et documents institutionnels}

\begin{apalist}
\item Organisation des Nations Unies pour l'éducation, la science et la culture.
      (1972). \emph{Convention concernant la protection du patrimoine mondial,
      culturel et naturel}. Paris: UNESCO.

\item Organisation des Nations Unies pour l'éducation, la science et la culture.
      (2021). \emph{Les musées dans le monde face à la pandémie de COVID-19}.
      Paris: UNESCO.
\end{apalist}

\section*{Documents électroniques et documentation technique}
\addcontentsline{toc}{section}{Documents électroniques et documentation technique}

\begin{apalist}
\item Chong, F., \& Carraro, G. (2006). \emph{Multi-tenant data architecture}.
      Retrieved from the Microsoft Developer Network web site:
      \url{https://learn.microsoft.com/en-us/previous-versions/dotnet/}

\item FinOps Foundation. (2023). \emph{FinOps framework}. Retrieved from
      \url{https://www.finops.org/framework/}

\item HashiCorp. (2024). \emph{Terraform documentation}. Retrieved from
      \url{https://developer.hashicorp.com/terraform/docs}

\item Khronos Group. (2021). \emph{glTF 2.0 specification}. Retrieved from
      \url{https://registry.khronos.org/glTF/specs/2.0/glTF-2.0.html}

\item Supabase. (2024). \emph{Row level security}. Retrieved from
      \url{https://supabase.com/docs/guides/database/postgres/row-level-security}

\item The Metropolitan Museum of Art. (2024). \emph{The Met collection API}.
      Retrieved from \url{https://metmuseum.github.io/}
\end{apalist}

\section*{Documents non publiés}
\addcontentsline{toc}{section}{Documents non publiés}

\begin{apalist}
\item Fondation Jean Félicien Gacha. (2026). \emph{Documents internes :
      inventaires, notices d'œuvres et plaquettes}. Document non publié,
      Bangoulap, Cameroun.

\item Institut Supérieur KEYCE Informatique \& Intelligence Artificielle.
      (2025). \emph{Guide de rédaction scientifique, version 3.0}. Document non
      publié, Bafoussam, Cameroun.
\end{apalist}
